\documentclass[aspectratio=169,10pt,spanish]{beamer}

\input{../../../_preambulo_beamer.tex}
\input{../../../_comandos_beamer.tex}

\selectlanguage{spanish}

\title{MA1001B - Modelación Estadística para la Toma de Decisiones}
\subtitle{Lección 07: Técnicas de conteo II}
\author{}
\institute{Tecnol\'ogico de Monterrey}
\date{}

\begin{document}

\begin{frame}[plain]
  \vspace{1.2cm}
  {\Large\bfseries MA1001B - Modelación Estadística para la Toma de Decisiones\par}
  \vspace{0.7cm}
  {\large Lección 07: Categorías repetidas y asignaciones\par}
  \vspace{0.7cm}
  {\color{TecRojo}\rule{\textwidth}{1.2pt}}
  \vspace{0.8cm}

  Facultad MA1001B\\[0.3cm]
  Periodo\\[0.3cm]
  Tecnol\'ogico de Monterrey
\end{frame}

\begin{frame}{La pregunta de asignación}
  \begin{alertblock}{Pregunta guía}
    ¿Cómo cambian las categorías repetidas un conteo, y cuándo un conteo
    correcto aún no determina una probabilidad?
  \end{alertblock}

  \vspace{0.4em}
  Al final de esta lección, usted puede:
  \begin{itemize}
    \item contar secuencias con etiquetas de categoría repetidas;
    \item interpretar coeficientes multinomiales como asignaciones de tamaño fijo;
    \item contar asignaciones no negativas con estrellas y barras;
    \item explicar por qué los perfiles posibles no tienen que ser igualmente
      probables.
  \end{itemize}

  \vfill
  \scriptsize Todas las solicitudes, colas, centros y mecanismos de enrutamiento son sintéticos.
\end{frame}

\begin{frame}{Las permutaciones con repetición eliminan duplicados}
  \small
  Ocho etiquetas de cola contienen tres Prioridad, tres Estándar y dos
  Auditoría. Intercambiar dos etiquetas iguales no crea una secuencia nueva.

  \[
    \frac{n!}{n_1!n_2!\cdots n_k!}
    =\frac{8!}{3!3!2!}=560
  \]

  \begin{alertblock}{El conteo ingenuo sobreestima el espacio muestral}
    \(8!=40{,}320\) trata las etiquetas repetidas como distinguibles. Cada
    secuencia distinta aparece \(3!3!2!=72\) veces en ese conteo.
  \end{alertblock}

  \textbf{Prueba de decisión:} ¿Qué intercambios cambian el resultado final?
\end{frame}

\begin{frame}[fragile]{Paso 1: La fórmula y la enumeración coinciden}
  \begin{columns}[T]
    \begin{column}{0.42\textwidth}
      \small
      \begin{lstlisting}
labels = (
    ["Prioridad"] * 3
    + ["Estandar"] * 3
    + ["Auditoria"] * 2
)
distinct = set(permutations(labels))
      \end{lstlisting}

      \begin{block}{Salida verificada}
        Conteo ingenuo: 40,320\\
        Fórmula: 560\\
        Secuencias enumeradas: 560\\
        Factor de duplicado: 72
      \end{block}
    \end{column}
    \begin{column}{0.58\textwidth}
      \centering
      \includegraphics[width=\textwidth]{../figuras/asignaciones_01_permutaciones_repetidas.png}
    \end{column}
  \end{columns}
\end{frame}

\begin{frame}{Paso 2: El mismo coeficiente asigna solicitudes etiquetadas}
  \begin{columns}[T]
    \begin{column}{0.40\textwidth}
      \small
      Asigne ocho solicitudes etiquetadas a tres colas etiquetadas con tamaños
      3, 3 y 2:
      \[
        \binom{8}{3,3,2}=\frac{8!}{3!3!2!}=560.
      \]

      En las 560 asignaciones, R01 aparece en:
      \begin{itemize}
        \item Prioridad: 210;
        \item Estándar: 210;
        \item Auditoría: 140.
      \end{itemize}

      La partición 3:3:2 sigue los tamaños fijos de cola.
    \end{column}
    \begin{column}{0.60\textwidth}
      \centering
      \includegraphics[width=\textwidth]{../figuras/asignaciones_02_categorias_multinomiales.png}
    \end{column}
  \end{columns}
\end{frame}

\begin{frame}{Paso 3: Contar perfiles de ocupación}
  \small
  Distribuya \(m\) unidades idénticas entre \(r\) categorías etiquetadas. Un
  perfil \((x_1,\ldots,x_r)\) cumple
  \[
    x_1+\cdots+x_r=m, \qquad x_i\geq0.
  \]

  El número de perfiles es
  \[
    \binom{m+r-1}{r-1}.
  \]

  Para ocho unidades y cuatro centros de servicio etiquetados:
  \[
    \binom{8+4-1}{4-1}=\binom{11}{3}=165.
  \]

  \begin{exampleblock}{Evidencia computacional verificada}
    La enumeración exhaustiva en Python también devuelve 165 perfiles de
    ocupación ordenados no negativos.
  \end{exampleblock}
\end{frame}

\begin{frame}{El límite: un perfil posible no es necesariamente equiprobable}
  \begin{columns}[T]
    \begin{column}{0.37\textwidth}
      \small
      Tratar los 165 perfiles como igualmente probables da
      \[
        \frac{1}{165}=0.00606061.
      \]

      El enrutamiento uniforme independiente da
      \[
        P(x)=\frac{8!}{x_1!\cdots x_4!}
        \left(\frac14\right)^8.
      \]

      Las estimaciones con semilla 42 de 500{,}000 simulaciones siguen de
      cerca estas probabilidades exactas desiguales.
    \end{column}
    \begin{column}{0.63\textwidth}
      \centering
      \includegraphics[width=\textwidth]{../figuras/asignaciones_03_estrellas_barras_limite_probabilidad.png}
    \end{column}
  \end{columns}
\end{frame}

\begin{frame}{Verificación de conceptos}
  \begin{enumerate}
    \item ¿Cuántas secuencias contienen tres etiquetas P, dos S y dos A?
    \item ¿Por qué el mismo coeficiente cuenta asignaciones de siete
      solicitudes etiquetadas a colas de tamaños 3, 2 y 2?
    \item ¿De cuántas formas se pueden distribuir ocho unidades idénticas de
      búfer entre tres centros etiquetados?
    \item ¿Son esos perfiles de ocupación igualmente probables cuando ocho
      solicitudes etiquetadas se enrutan de forma independiente? Explique.
  \end{enumerate}

  \vspace{0.4em}
  \begin{block}{Conexión con la Lección 08}
    La siguiente lección formaliza las reglas de probabilidad y las tablas de
    contingencia. Un conteo se vuelve probabilidad solo después de justificar
    el modelo de probabilidad.
  \end{block}

  \vfill
  \scriptsize Fundamento de probabilidad: OpenStax,
  \textit{Introductory Business Statistics 2e}, Capítulo 3, Sección 3.1.
  Las fórmulas combinatorias son una extensión del curso. CC BY-NC-SA 4.0.
\end{frame}

\appendix

\begin{frame}{Verificación de conceptos: respuestas sugeridas}
  \begin{enumerate}
    \item $\frac{7!}{3!2!2!}=210$ secuencias distintas.
    \item Elegir las tres solicitudes P y luego dos del resto para S deja
      dos para A: $\binom{7}{3}\binom{4}{2}=210$. Esto iguala el
      coeficiente multinomial $\binom{7}{3,2,2}$.
    \item Estrellas y barras da
      $\binom{8+3-1}{3-1}=\binom{10}{2}=45$ perfiles.
    \item No. Las secuencias de enrutamiento etiquetadas se colapsan en
      perfiles de ocupación con distintas multiplicidades, de modo que los
      perfiles en general tienen distintas probabilidades.
  \end{enumerate}

  \begin{block}{Nota para la persona instructora}
    Exija definir el resultado final e identificar cuáles objetos y
    categorías están etiquetados antes de aceptar una fórmula.
  \end{block}
\end{frame}

\end{document}
